% Part: first-order-logic 
% Chapter: axiomatic-deduction 
% Section: provability

% verification of properties of provability needed for maximally 
% consistent sets in the completeness chapter.

\documentclass[../../../include/open-logic-section]{subfiles}

\begin{document}

\iftag{FOL}
      {\olfileid{fol}{axd}{prv}}
      {\olfileid{pl}{axd}{prv}}

\olsection{Sipat-sipat \usetoken{S}{derivability}}

\begin{prop}[Monotonisitas] Yèn $\Gamma \subseteq \Delta$ lan $\Gamma \Proves !A$,
mula uga $\Delta \Proves !A$. \end{prop}

\begin{proof} Saben $\Gamma_0 \subseteq \Gamma$ kang winates uga
himpunan pérangan winates saka~$\Delta$, mula !!a{derivation}
kanggo $\Gamma_0 \Proves !A$ uga nuduhaké $\Delta \Proves !A$.
\end{proof}

\begin{prop}\ollabel{prop:provability} $\Gamma \Proves !A$ yèn lan
namung yèn $\Gamma\cup \{\lnot!A \}$ ora konsisten. \end{prop}

\begin{prop} \begin{enumerate}
\item \ollabel{prop:provability-contr} Yèn $\Gamma \Proves !A$ lan $\Gamma
\cup \{ !A\} \Proves \lfalse$, mula $\Gamma$ ora konsisten.

\item \ollabel{prop:provability-lnot} Yèn $\Gamma \cup \{!A\}
\Proves \lfalse$, mula $\Gamma \Proves \lnot !A$.

\item \ollabel{prop:provability-exhaustive} Yèn $\Gamma \cup \{!A\}
\Proves \lfalse$ lan $\Gamma \cup \{\lnot !A\} \Proves
\lfalse$, mula $\Gamma \Proves \lfalse$.

\tagitem{prvOr}{\ollabel{prop:provability-lor-left} Yèn $\Gamma \cup \{!A\}
\Proves \lfalse$ lan $\Gamma \cup \{!B\} \Proves
\lfalse$, mula $\Gamma \cup \{!A \lor !B\} \Proves \lfalse$.}{}

\tagitem{prvOr}{\ollabel{prop:provability-lor-right} Yèn $\Gamma
\Proves !A$ utawa $\Gamma \Proves !B$, mula $\Gamma
\Proves !A \lor !B$.}{}

\tagitem{prvAnd}{\ollabel{prop:provability-land-left} Yèn $\Gamma
\Proves !A \land !B$, mula $\Gamma \Proves !A$ lan
$\Gamma \Proves !B$.}{}

\tagitem{prvAnd}{\ollabel{prop:provability-land-right} Yèn $\Gamma
\Proves !A$ lan $\Gamma \Proves !B$, mula $\Gamma
\Proves !A \land !B$.}{}

\tagitem{prvIf}{\ollabel{prop:provability-mp} Yèn $\Gamma \Proves
!A$ lan $\Gamma \Proves !A \lif !B$, mula $\Gamma
\Proves !B$.}{}

\tagitem{prvIf}{\ollabel{prop:provability-lif} Yèn $\Gamma \Proves
\lnot !A$ utawa $\Gamma \Proves !B$, mula $\Gamma \Proves
!A \lif !B$.}{} \end{enumerate} \end{prop}

\begin{proof} 
\begin{enumerate}
\item Anggepen $\Gamma_0 \cup \{!A\} \Proves \lfalse$.

\begin{derivation} 
1. & $\Gamma_0 \cup \{!A\} \Proves \lfalse$ & asumsi \\
2. & $\Gamma_1 \Proves !A$ & asumsi\\ 
3. & $\Gamma_0 \cup \Gamma_1 \Proves \lfalse$ & cut \\ 
\end{derivation}

Amarga $\Gamma_0 \subseteq \Gamma$ lan $\Gamma_1 \subseteq \Gamma$,
$\Gamma_0 \cup \Gamma_1 \subseteq \Gamma$, mula $\Gamma \Proves \lfalse$.

\item Anggepen $\Gamma_0 \cup\{!A\} \Proves \lfalse$.

\begin{derivation} 
1. & $\Gamma_0 \cup\{!A\} \Proves \lfalse$ & asumsi \\ 
2. & $\Gamma_0 \Proves !A \lif \lfalse$ & Téoréma Deduksi\\
3. & $\Gamma_0 \Proves (!A \lif \lfalse) \lif \lnot !A$ & Ax0\\ 
4. & $\Gamma_0 \Proves \lnot !A$ & MP 2, 3 \\ 
\end{derivation}

\item Anggepen $\Gamma_0 \cup \{ !A \} \Proves \lfalse$ lan
$\Gamma_1 \cup \{\lnot !A \} \Proves \lfalse$.

\begin{derivation} 
1. & $\Gamma_0 \cup\{!A\} \Proves \lfalse$ & asumsi \\
2. & $\Gamma_0 \Proves !A \lif \lfalse$ & Téoréma Deduksi \\ 
3. & $\Gamma_1 \cup \{\lnot !A \} \Proves \lfalse$ & asumsi\\ 
4. & $\Gamma_0 \Proves (!A \lif \lfalse) \lif \lnot !A$ & Ax0\\ 
5. & $\Gamma_0 \Proves \lnot !A$ & MP 2, 4\\ 
6. & $\Gamma_0 \cup \Gamma_1 \Proves \lfalse$ & Cut 3, 5 \\
\end{derivation}

Amarga $\Gamma_0 \subseteq \Gamma$ lan $\Gamma_1 \subseteq \Gamma$,
$\Gamma_0 \cup \Gamma_1 \subseteq \Gamma$. Mula $\Gamma \Proves \lfalse$.

% prop:provability-lor-left 
\tagitem{defOr}{}{
\iftag{probOr}{Gladhen.}{Gladhen.}}

% prop:provability-lor-right 
\tagitem{defOr}{}{ \iftag{probOr}{Gladhen.}{
Anggepen $\Gamma_0 \Proves !A$.

\begin{derivation} 
1. & $\Gamma_0 \Proves !A$ & asumsi \\ 
2. & $\Gamma_0 \Proves !A \lif (!A \lor !B)$ & Ax0\\ 
3. & $\Gamma_0 \Proves !A \lor !B$ & MP 1, 2 \\
\end{derivation}

Mula $\Gamma \Proves !A \lor !B$. Bukti kanggo kahanan
$\Gamma \Proves !B$ padha carané.}}

% prop:provability-land-left 
\tagitem{defAnd}{}{ \iftag{probAnd}{Gladhen}{
Anggepen $\Gamma_0 \Proves !A \land !B$

\begin{derivation} 
1. & $\Gamma_0 \Proves !A \land !B$ & asumsi \\ 
2. & $\Gamma_0 \Proves (!A\land !B) \lif !A $ & Ax0\\ 
3. & $\Gamma_0 \Proves !A$ & MP 1, 2 \\ 
\end{derivation}

Mula $\Gamma \Proves !A$. !!{derivation} kang padha carané
uga nuduhaké yèn $\Gamma \Proves !B$.}}

% prop:provability-land-right
\tagitem{defAnd}{}{\iftag{probAnd}{Gladhen.}{Gladhen.}}

% prop:provability-mp 
\tagitem{defIf}{}{ \iftag{probIf}{Gladhen.}{
Anggepen $\Gamma_0 \Proves !A$ lan $\Gamma_1 \Proves !A \lif !B $.

\begin{derivation} 
1. & $\Gamma_0 \Proves !A$ & asumsi \\ 
2. & $\Gamma_1 \Proves !A \lif !B $ & asumsi\\ 
3. & $\Gamma_0 \cup \Gamma_1 \Proves !B$ & MP 1, 2\\
\end{derivation}

Amarga $\Gamma_0 \cup \Gamma_1 \subseteq \Gamma$, iki ateges $\Gamma
\Proves !B$.}}

% prop:provability-lif 
\tagitem{defIf}{}{\iftag{probIf}{Gladhen.}{Dhisik,
anggepen $\Gamma \Proves \lnot !A$.
    
\begin{derivation} 
1. & $\Gamma_0 \Proves \lnot !A$ & asumsi \\ 
2. & $\Gamma_0 \Proves \lnot !A \lif (!A \lif !B) $ & Ax0\\ 
3. & $\Gamma_0 \Proves !A \lif !B$ & MP 1,
2\\ \end{derivation}

Bukti kanggo $\Gamma \Proves !B$ padha carané:

\begin{derivation} 
1. & $\Gamma_0 \Proves !B$ & asumsi \\ 
2. & $\Gamma_0 \Proves !B \lif (!A \lif !B) $ & Ax0\\ 
3. & $\Gamma_0 \Proves !A \lif !B$ & MP 1, 2\\
\end{derivation}}}

\end{enumerate} 
\end{proof}

\begin{probtag}{probOr,probAnd,probIf} Rampungna bukti kanggo
\olref[fol][axd][prv]{prop:provability}. 
\end{probtag}

\begin{thm}[Generalisasi] \ollabel{thm:Generalization} Yèn $\Gamma \Proves
!A$ lan $x$ ora bebas ing !!{formula} endi waé sajroning $\Gamma$,
mula $\Gamma \Proves \lforall[x][!A]$.
\end{thm}

\begin{proof} Kanthi induksi tumrap $\mathsf{Thm}_\Gamma$: yèn $!A$
iku aksioma, $\lforall[x][!A]$ uga aksioma. Yèn $!A\in \Gamma$,
$x$ ora bebas ing $!A$, mula $!A \lif \lforall[x][!A]$
iku aksioma (\textbf{Ax3}), lan kanthi MP kita olèh $\Gamma
\Proves \lforall[x][!A]$. Saiki anggepen $!A$ diturunké kanthi MP
amarga $\Gamma \Proves !B$ lan $\Gamma \Proves !B \lif !A$.
Miturut hipotesis induksi, $\Gamma \Proves \lforall[x][!B]$
lan $\Gamma \Proves \lforall[x][( !B \lif !A)]$.
Nanging miturut \textbf{Ax2} uga \[ \Gamma \Proves
\lforall[x][( !B \lif !A)] \lif (\lforall[x][ !B] \lif \lforall[x][!A]), \]
mula rong panganggoné MP ngasilaké $\Gamma \Proves
\lforall[x][!A]$ kaya kang dikarepaké. \end{proof}

\begin{thm}\ollabel{thm:WeakGen} (\emph{Generalisasi Lemah tumrap Konstanta})
Yèn $\Gamma \Proves !A$ lan $c$ iku !!a{constant} kang ora
katon ing $\Gamma$, mula ana variabel $x$ kang ora katon
ing $!A$ saéngga $\Gamma
\Proves \lforall[x][\Subst{!A}{x}{c}]$, lan buktiné ora nganggo $c$.
\end{thm}

\begin{proof} Anggep $!A_1,\ldots,!A_n$ minangka bukti $!A$ saka
$\Gamma$, saéngga $!A=!A_n$. Pilih variabel $x$ kang ora katon
ing $!A_1,\ldots,!A_n$, banjur delengen runtutan anyar:
$\Subst{!A_1}{x}{c},\ldots,\Subst{!A_n}{x}{c}.$ Runtutan iki
minangka bukti $\Subst{!A}{x}{c}$ saka $\Gamma$. Kanggo saben
$i=1,\ldots,n$: \begin{itemize} \item yèn $!A_i$ iku aksioma,
$\Subst{!A_i}{x}{c}$ uga aksioma; \item yèn $!A_i \in \Gamma$,
amarga $c$ ora katon ing $\Gamma$, kita olèh
$\Subst{!A_i}{x}{c} = !A_i \in \Gamma$; \item yèn $!A_i$
diturunké saka $!A_j$ lan $!A_j \lif !A_i$, mula
$\Subst{!A_i}{x}{c}$ diturunké kanthi MP saka
$\Subst{!A_j}{x}{c}$ lan $\Subst{(!A_j \lif !A_i)}{x}{c}
= \Subst{!A_j}{x}{c} \lif \Subst{!A_i}{x}{c}$.
\end{itemize} Cetha yèn !!{constant} $c$ ora katon manèh
ing runtutan anyar iki. Saiki anggep $\Gamma'$ ngemot !!{formula}s
saka $\Gamma$ kang katon ing !!{derivation}
$\Subst{!A}{x}{c}$; mula $x$ ora bebas ing !!{formula} endi waé
sajroning $\Gamma'$ lan $\Gamma' \Proves \Subst{!A}{x}{c}$.
Miturut generalisasi (\olref{thm:Generalization}), $\Gamma'
\Proves \lforall[x][\Subst{!A}{x}{c}]$, lan miturut
monotonisitas $\Gamma \Proves
\lforall[x][\Subst{!A}{x}{c}]$, kaya kang dikarepaké.
\end{proof}

\begin{explain}
Tujuan kita yaiku ngganti syarat ing \olref{thm:WeakGen} yèn $x$
ora katon babar pisan ing $!A$ nganggo syarat kang luwih lemah:
$x$ ora bebas ing $!A$ lan !!{free for} $c$ ing $!A$.
Iki bisa ditindakake kanthi ngganti !!{variable} kaiket;
mula iki kang bakal dibuktèkaké luwih dhisik.
\end{explain}

\begin{lem}
\ollabel{lem:free-for} 
Yèn $x$ bebas kanggo $c$ ing $!A$ lan $y$ ora bebas ing $!A$,
mula $x$ iku !!{free for} $y$ ing $\Subst{!A}{y}{c}$.
\end{lem}

\begin{proof} Yèn $x$ ora !!{free for} $y$ ing
$\Subst{!A}{y}{c}$, mula ana kemunculan bebas saka $y$ ing
$\Subst{!A}{y}{c}$ kang mapan ing cakupan kuantor
$\lforall[x]$. Nanging, miturut hipotesis, $y$ ora bebas
ing $!A$, mula saben kemunculan kaya mangkono asalé saka
substitusi $\Subst{}{y}{c}$. Dadi ana kemunculan $c$
ing cakupan $\lforall[x]$, saéngga $x$ ora
!!{free for} $c$ ing $!A$.
\end{proof}

\begin{lem}[Panggantèn Variabel Kaiket]
\ollabel{lem:ChangeBdVar} 
Yèn $x$ lan $y$ ora bebas ing $!A$ lan kaloroné
!!{free for} $c$ ing $!A$, mula
$\Proves \lforall[x][\Subst{!A}{x}{c}] \equiv
\lforall[y][\Subst{!A}{y}{c}]$ (lan buktiné ora nganggo $c$).
\end{lem}

\begin{proof} 
Miturut hipotesis, $y$ iku !!{free for} $c$ ing $!A$ lan $x$
ora bebas ing $!A$. Mula, miturut \olref{lem:free-for}, $y$
iku !!{free for} $x$ ing $\Subst{!A}{x}{c}$. Saka kéné,
$\lforall[x][\Subst{!A}{x}{c} \lif \Subst{!A}{x}{c} \Subst{}{y}{x}]$
iku aksioma (\textbf{Ax1}). Amarga $x$ ora bebas ing $!A$,
kita olèh $\Subst{!A}{x}{c} \Subst{}{y}{x} = \Subst{!A}{y}{c}$,
mula
\[
\Proves \lforall[x][\Subst{!A}{x}{c}] \lif \Subst{!A}{y}{c}, 
\] 
lan miturut Téoréma Deduksi,
$\lforall[x][\Subst{!A}{x}{c}] \Proves \Subst{!A}{y}{c}$.
Amarga $y$ ora bebas ing $!A$, dhèwèké uga ora bebas ing
$\lforall[x][\Subst{!A}{x}{c}]$, mula miturut Generalisasi,
$\lforall[x][\Subst{!A}{x}{c}] \Proves \lforall[y][\Subst{!A}{y}{c}]$.
Téoréma Deduksi sapisan manèh mènèhi $\Proves
\lforall[x][\Subst{!A}{x}{c}] \lif \lforall[y][\Subst{!A}{y}{c}]$.
Bukti kanggo $\Proves \lforall[y][\Subst{!A}{y}{c}] \lif \lforall[x]
[\Subst{!A}{x}{c}]$ simètris, mula dudutané manut aturan Taut.
\end{proof}

\begin{thm}[Generalisasi Kuwat tumrap Konstanta]
\ollabel{thm:StrongGen} 
Yèn $\Gamma \Proves !A$, !!{constant} $c$ ora katon ing
$\Gamma$, lan $x$ ora bebas ing $!A$ nanging !!{free for}
$c$ ing $!A$, mula $\Gamma \Proves
\lforall[x][\Subst{!A}{x}{c}]$.
\end{thm}

\begin{proof} 
Amarga $\Gamma \Proves !A$ lan $c$ ora katon ing $\Gamma$,
miturut Generalisasi Lemah ana !!a{variable} $y$ kang ora
katon ing $!A$ saéngga
$\Gamma \Proves \lforall[y][\Subst{!A}{y}{c}]$.
Amarga $y$ ora katon babar pisan ing $!A$, dhèwèké ora bebas
ing $!A$ lan uga !!{free for} $c$ ing $!A$.
Miturut hipotesis, $x$ uga ora bebas ing $!A$ lan !!{free for}
$c$ ing $!A$. Mula syarat kanggo ngganti !!{variable} kaiket
wis kacukupan, saéngga
$\Proves \lforall[x][\Subst{!A}{x}{c}] \equiv
\lforall[y][\Subst{!A}{y}{c}]$, lan akibaté $\Gamma \Proves
\lforall[x][\Subst{!A}{x}{c}]$.
\end{proof}

\begin{thm} 
\ollabel{thm:provability-quantifiers}
\begin{tagenumerate}{prvEx,prvAll} 
\tagitem{prvEx}{Yèn $\Gamma \Proves !A(t)$, mula $\Gamma \Proves
  \lexists[x][!A(x)]$.}{} 
\tagitem{prvAll}{Yèn $\Gamma \Proves \lforall[x][!A(x)]$, mula $\Gamma
  \Proves !A(t)$.}{}
\end{tagenumerate} 
\end{thm}

\begin{proof} 
\begin{tagenumerate}{prvEx,prvAll} 
%%Need axioms for exists
\tagitem{prvEx}{Gladhen.}{}
\tagitem{prvAll} {Anggepen $\Gamma_0 \Proves \lforall[x][!A(x)]$.

\begin{derivation} 
1. & $\Gamma_0 \Proves \lforall[x][!A(x)]$ & asumsi \\
2. & $\Gamma_0 \Proves \lforall[x][!A(x)] \lif \Subst{!A}{t}{x}$ & Ax1 \\ 
3. & $\Gamma_0 \Proves \Subst{!A}{t}{x}$ & MP 1, 2 \\ 
\end{derivation}}{}

\end{tagenumerate} 
\end{proof}

\end{document}
